\documentclass[notes=show]{beamer}
%%%%%%%%%%%%%%%%%%%%%%%%%%%%%%%%%%%%%%%%%%%%%%%%%%%%%%%%%%%%%%%%%%%%%%%%%%%%%%%%%%%%%%%%%%%%%%%%%%%%%%%%%%%%%%%%%%%%%%%%%%%%%%%%%%%%%%%%%%%%%%%%%%%%%%%%%%%%%%%%%%%%%%%%%%%%%%%%%%%%%%%%%%%%%%%%%%%%%%%%%%%%%%%%%%%%%%%%%%%%%%%%%%%%%%%%%%%%%%%%%%%%%%%%%%%%
\AtBeginSubsection[
  {\frame<beamer>{\frametitle{Outline}   
    \tableofcontents[currentsection,currentsubsection]}}%
]%
{
  \frame<beamer>{ 
    \frametitle{Outline}   
    \tableofcontents[currentsection,currentsubsection]}
}
\usepackage{ulem}
\usepackage{pdflscape}
\setbeamertemplate{caption}{\insertcaption}
\usepackage{color,xcolor,ucs}
\usepackage{beamerthemesplit}
\usepackage{graphicx}
\usepackage{amsmath}
\usepackage{adjustbox}
\usepackage{pdfpages} 
\definecolor{mintgreen}{rgb}{0.6, 1.0, 0.6}
\makeatletter
\@addtoreset{subfigure}{framenumber}% subfigure counter resets every frame
\makeatother
\newcommand*{\Scale}[2][4]{\scalebox{#1}{$#2$}}%
\newcommand*{\Resize}[2]{\resizebox{#1}{!}{$#2$}}%
\usepackage{amsfonts}
\usepackage{booktabs}
\usepackage{color,soul}
\usepackage{pdfpages}
\setboolean{@twoside}{false}
\usepackage{amsmath}
\usepackage{pdflscape} 
\usepackage{eurosym}
\usepackage{amstext} % for \text
\DeclareRobustCommand{\officialeuro}{%
  \ifmmode\expandafter\text\fi
  {\fontencoding{U}\fontfamily{eurosym}\selectfont e}}
\usepackage[caption = false]{subfig}
\usepackage{appendixnumberbeamer} 
\usepackage{graphicx}
\usepackage{mathpazo}
\usepackage{hyperref}
\usepackage{multimedia}
\usepackage{graphicx}
\usepackage{multirow}
\usepackage{subfig}
\usepackage{verbatim}
\usepackage{booktabs, multicol, multirow}
\setcounter{MaxMatrixCols}{10}
\input{Macros.tex}
\AtBeginSection[]
{
  \begin{frame}[noframenumbering]
    \frametitle{Outline for section \thesection}
    \tableofcontents[currentsection]
  \end{frame}
}


\usepackage{xcolor,soul}
\renewcommand<>{\hl}[1]{\only#2{\beameroriginal{\hl}}{#1}}

% https://tex.stackexchange.com/questions/41683/why-is-it-that-coloring-in-soul-in-beamer-is-not-visible
\makeatletter
\newcommand\SoulColor{%
  \let\set@color\beamerorig@set@color
  \let\reset@color\beamerorig@reset@color}
\makeatother
\SoulColor
\AtBeginSection{\frame{\sectionpage}}

\newsavebox\mybox
\newenvironment{aquote}[1]
  {\savebox\mybox{#1}\begin{quote}\openautoquote\hspace*{-.7ex}}
  {\unskip\closeautoquote\vspace*{1mm}\signed{\usebox\mybox}\end{quote}}

\newtheorem{proposition}[theorem]{Proposition}
\newtheorem{Assumption}[theorem]{Assumption}
\newenvironment{stepenumerate}{\begin{enumerate}[<+->]}{\end{enumerate}}
\newenvironment{stepitemize}{\begin{itemize}[<+->]}{\end{itemize} }
\newenvironment{stepenumeratewithalert}{\begin{enumerate}[<+-| alert@+>]}{\end{enumerate}}
\newenvironment{stepitemizewithalert}{\begin{itemize}[<+-| alert@+>]}{\end{itemize} }
\usetheme{Madrid}
\makeatletter
\setbeamertemplate{footline}
{
  \leavevmode%
  \hbox{%
  \begin{beamercolorbox}[wd=.333333\paperwidth,ht=2.25ex,dp=1ex,center]{author in head/foot}%
    \usebeamerfont{author in head/foot}\insertshortauthor~~\beamer@ifempty{\insertshortinstitute}{}{(\insertshortinstitute)}
  \end{beamercolorbox}%
  \begin{beamercolorbox}[wd=.333333\paperwidth,ht=2.25ex,dp=1ex,center]{title in head/foot}%
    \usebeamerfont{title in head/foot}\insertshorttitle
  \end{beamercolorbox}%
  \begin{beamercolorbox}[wd=.333333\paperwidth,ht=2.25ex,dp=1ex,right]{date in head/foot}%
    \usebeamerfont{date in head/foot}\insertshortdate{}\hspace*{2em}
%    \insertframenumber{} / \inserttotalframenumber\hspace*{2ex} % DELETED
  \end{beamercolorbox}}%
  \vskip0pt%
}
\makeatother
\usepackage{pdfpages}
\begin{document}
	\title[]{\large Inference
	\vspace{0.5cm}
	}
	\author[Raffaele Saggio - Econ 628]{Raffaele Saggio}
	\institute[]{UBC}
	\date{\today \\
}
	\maketitle

\section{Statistics Boot Camp}

\begin{frame}{Statistical Properties of OLS}

\begin{itemize}
\item We've talked about the population OLS projection as an approximation
to the CEF {\small{}\medskip{}
}{\small\par}
\item This didn't require any assumptions about the data-generating process{\small{}\medskip{}
}{\small\par}
\item Next, we'll talk about the statistical properties of the OLS \textbf{estimator}
based on a sample of data, and how these depend on the DGP
\end{itemize}
\end{frame}

\begin{frame}{The OLS Estimator}

\begin{itemize}
\item Let's again work with the projection model:
\[
Y_{i}=X_{i}^{\prime}\beta+U_{i}\mbox{, \ensuremath{E\left[X_{i}U_{i}\right]=0}}
\]
\item We have a sample of $N$ individuals from the population{\small{}\medskip{}
}{\small\par}
\item {\small{}The OLS estimator of $\beta$ is}{\small\par}
\end{itemize}
\begin{center}
{\small{}$\hat{\beta}=\left(\tfrac{1}{N}\sum_{i}X_{i}X_{i}^{\prime}\right)^{-1}\left(\tfrac{1}{N}\sum_{i}X_{i}Y_{i}\right)$}{\small\par}
\par\end{center}

\begin{center}
$=\left(\tfrac{1}{N}\sum_{i}X_{i}X_{i}^{\prime}\right)^{-1}\left(\tfrac{1}{N}\sum_{i}X_{i}(X_{i}^{\prime}\beta+U_{i})\right)$
\par\end{center}

\begin{center}
$=\beta+\left(\tfrac{1}{N}\sum_{i}X_{i}X_{i}^{\prime}\right)^{-1}\left(\tfrac{1}{N}\sum_{i}X_{i}U_{i}\right)$
\par\end{center}

\end{frame}

\begin{frame}{OLS in Finite Samples}

\begin{itemize}
\item {\small{}First consider the finite-sample properties of $\hat{\beta}$}{\small\par}
\item Its expectation is
\end{itemize}
\begin{center}
$E\left[\hat{\beta}\right]=\beta+E\left[\left(\tfrac{1}{N}\sum_{i}X_{i}X_{i}^{\prime}\right)^{-1}\left(\tfrac{1}{N}\sum_{i}X_{i}U_{i}\right)\right]$
\par\end{center}

\begin{center}
$\neq\beta+\left(\tfrac{1}{N}\sum_{i}E\left[X_{i}X_{i}^{\prime}\right]\right)^{-1}\left(\tfrac{1}{N}\sum_{i}E\left[X_{i}U_{i}\right]\right)$
\par\end{center}

\begin{center}
$=\beta$
\par\end{center}
\begin{itemize}
\item In general $\hat{\beta}$ is a biased estimate of the population projection
coefficient vector $\beta${\small{}\medskip{}
}{\small\par}
\item An exception is when the CEF is linear so $E\left[Y_{i}|X_{i}\right]=X_{i}^{\prime}\beta$
(just apply the law of iterated expectations)
\end{itemize}
\end{frame}

\begin{frame}{Fixed vs. Stochastic Regressors}

\begin{itemize}
\item {\small{}You may be surprised to learn that OLS is biased in finite
samples for the population projection\medskip{}
}{\small\par}
\item {\small{}The properties of OLS are traditionally derived assuming
that the regressor matrix $X$ is non-stochastic, i.e., fixed in repeated
samples\medskip{}
}{\small\par}
\item {\small{}Sometimes (e.g. in an experiment) the researcher gets to
control the distribution of the regressors\medskip{}
}{\small\par}
\item {\small{}When $X$ is non-stochastic we have}{\small\par}
\end{itemize}
\begin{center}
{\small{}$E\left[\hat{\beta}\right]=\beta+E\left[\left(\tfrac{1}{N}\sum_{i}X_{i}X_{i}^{\prime}\right)^{-1}\left(\tfrac{1}{N}\sum_{i}X_{i}U_{i}\right)\right]$}{\small\par}
\par\end{center}

\begin{center}
{\small{}$=\beta+\left(\tfrac{1}{N}\sum_{i}X_{i}X_{i}^{\prime}\right)^{-1}\left(\tfrac{1}{N}\sum_{i}E\left[X_{i}U_{i}\right]\right)$}{\small\par}
\par\end{center}

\begin{center}
{\small{}$=\beta$}{\small\par}
\par\end{center}

\end{frame}

\begin{frame}{Fixed vs. Stochastic Regressors}

\begin{itemize}
\item Whether it makes sense to consider the regressors fixed or random
depends on the context{\small{}\medskip{}
}{\small\par}
\item Do we want a best approximation to the CEF using the distribution
of $X_{i}$ in our sample, or in the population?

\begin{itemize}
\item {\small{}Note: sometimes your sample is ``the population'' (e.g.
everyone in Iowa)\medskip{}
}{\small\par}
\end{itemize}
\item For now let's consider the traditional fixed-regressor case
\end{itemize}
\end{frame}

\begin{frame}{The Variance of OLS}

\begin{itemize}
\item {\small{}In matrix form the variance of the OLS estimator is}{\small\par}
\end{itemize}
\begin{center}
{\small{}$E\left[\left(\hat{\beta}-E\left[\hat{\beta}\right]\right)\left(\hat{\beta}-E\left[\hat{\beta}\right]\right)^{\prime}\right]=E\left[\left(\hat{\beta}-\beta\right)\left(\hat{\beta}-\beta\right)^{\prime}\right]$}{\small\par}
\par\end{center}

\begin{center}
{\small{}$=E\left[\left(X^{\prime}X\right)^{-1}X^{\prime}UU^{\prime}X\left(X^{\prime}X\right)^{-1}\right]$}{\small\par}
\par\end{center}

\begin{center}
{\small{}$=\left(X^{\prime}X\right)^{-1}X^{\prime}E\left[UU^{\prime}\right]X\left(X^{\prime}X\right)^{-1}$}{\small\par}
\par\end{center}

\begin{center}
{\small{}$=\left(X^{\prime}X\right)^{-1}X^{\prime}\Omega X\left(X^{\prime}X\right)^{-1}$}{\small\par}
\par\end{center}
\begin{itemize}
\item {\small{}$\Omega$ is the $N\times N$ symmetric covariance matrix
with $ij$-th element $Cov(U_{i},U_{j})$}{\small\par}
\end{itemize}
\end{frame}

\begin{frame}{The Variance of OLS}

\begin{itemize}
\item {\small{}The ``}meat{\small{}'' in the middle of our sandwich covariance
matrix can be written}{\small\par}
\end{itemize}
\begin{center}
{\small{}$X^{\prime}\Omega X=\sum_{i}X_{i}X_{i}^{\prime}\sigma_{i}^{2}+2\sum_{i<j}X_{i}X_{j}^{\prime}\sigma_{ij}$}{\small\par}
\par\end{center}
\begin{itemize}
\item {\small{}Here $\sigma_{i}^{2}=Var\left(U_{i}\right)$ and $\sigma_{ij}=Cov\left(U_{i},U_{j}\right)$\medskip{}
}{\small\par}
\item {\small{}There are $N(N+1)/2$ unknown terms in this expression\medskip{}
}{\small\par}
\item {\small{}Estimating all of these terms is intractable without further
assumptions}{\small\par}
\end{itemize}
\end{frame}

\begin{frame}{The Variance of OLS}

\begin{itemize}
\item {\small{}Suppose the regression errors are independent and homoskedastic,
so $Cov\left(U_{i},U_{j}\right)=0$ and $Var\left(U_{i}\right)=\sigma^{2}\ \forall i$.
Then}{\small\par}
\end{itemize}
\begin{center}
{\small{}$Var\left(\hat{\beta}\right)=\left(X^{\prime}X\right)^{-1}X^{\prime}\left(\sigma^{2}I\right)X\left(X^{\prime}X\right)^{-1}$}{\small\par}
\par\end{center}

\begin{center}
{\small{}$=\sigma^{2}\left(X^{\prime}X\right)^{-1}$}{\small\par}
\par\end{center}
\begin{itemize}
\item {\small{}The conventional estimator of $\sigma^{2}$ is}

\begin{center}
{\small{}$\hat{\sigma}^{2}=\dfrac{1}{N-K}\sum_{i}\hat{U}_{i}^{2}$}{\small\par}
\par\end{center}

\end{itemize}
\end{frame}

\begin{frame}{Degrees of Freedom Correction}

\begin{center}
{\footnotesize{}$\hat{\sigma}^{2}=\dfrac{1}{N-K}\sum_{i}\hat{U}_{i}^{2}$}{\footnotesize\par}
\par\end{center}
\begin{itemize}
\item {\footnotesize{}Why $1/(N-K)$ instead of $1/N$? \medskip{}
}{\footnotesize\par}
\item {\footnotesize{}This is a degrees of freedom correction to account
for finite-sample bias\medskip{}
}{\footnotesize\par}
\item {\footnotesize{}$\hat{U}=M_{X}Y$, so}{\footnotesize\par}
\end{itemize}
\begin{center}
{\footnotesize{}$Var\left(\hat{U}\right)=E\left[M_{X}YY^{\prime}M_{X}\right]$}{\footnotesize\par}
\par\end{center}

\begin{center}
{\footnotesize{}$=M_{X}E\left[UU^{\prime}\right]M_{X}$}{\footnotesize\par}
\par\end{center}

\begin{center}
{\footnotesize{}$=\sigma^{2}M_{X}$.}{\footnotesize\par}
\par\end{center}
\begin{itemize}
\item {\footnotesize{}Then $Var\left(\hat{U}_{i}\right)=\sigma^{2}(1-P_{ii})$,
where $P_{ii}$ is the $i$'th diagonal entry in $P_{X}$ }{\footnotesize\par}
\end{itemize}
\end{frame}

\begin{frame}{Degrees of Freedom Correction}

\begin{center}
{\footnotesize{}$Var\left(\hat{U}_{i}\right)=\sigma^{2}(1-P_{ii})$}{\footnotesize\par}
\par\end{center}
\begin{itemize}
\item {\footnotesize{}It turns out that $\sum_{i}P_{ii}=K$:}{\footnotesize\par}
\end{itemize}
\begin{center}
{\footnotesize{}$\sum_{i}P_{ii}=tr\left(X\left(X^{\prime}X\right)^{-1}X^{\prime}\right)=tr\left(X^{\prime}X\left(X^{\prime}X\right)^{-1}\right)=tr(I)=K$}{\footnotesize\par}
\par\end{center}
\begin{itemize}
\item {\footnotesize{}This means that the expectation of the mean squared
residual is}{\footnotesize\par}
\end{itemize}
\begin{center}
{\footnotesize{}$E\left[\frac{1}{N}\sum_{i}\hat{U}_{i}^{2}\right]=\frac{1}{N}\sum_{i}E\left[\hat{U}_{i}^{2}\right]=\frac{1}{N}\sum\sigma^{2}(1-P_{ii})=\sigma^{2}\cdot\frac{N-K}{N}$}{\footnotesize\par}
\par\end{center}
\begin{itemize}
\item {\footnotesize{}We therefore need to multiply by $(N/(N-K))$ to get
an unbiased estimate of $\sigma^{2}$ (just like we divide by $N-1$
instead of $N$ to get unbiased variance estimates)}{\footnotesize\par}
\end{itemize}
\end{frame}

\begin{frame}{Finite-sample Inference}

\begin{itemize}
\item {\small{}We've shown how to get an unbiased estimate of $\hat{\sigma}^{2}$,
and therefore of $Var\left(\hat{\beta}\right)$, under homoskedasticity\medskip{}
}{\small\par}
\item {\small{}But this isn't enough to allow us to test hypotheses or form
confidence intervals}{\small\par}
\end{itemize}
{\small{}\medskip{}
}{\small\par}
\begin{itemize}
\item {\small{}To conduct inference we need to know the distribution of
$\hat{\beta}$\medskip{}
}{\small\par}
\item {\small{}The traditional approach is to assume that the $U_{i}$ are
$iid$ and normally distributed, which implies that $\hat{\beta}$
is normal and $(\hat{\beta}_{k}-\beta_{k})/\hat{Var}\left(\hat{\beta}_{k}\right)$
follows a $t$-distribution with $(N-K)$ degrees of freedom\medskip{}
}{\small\par}
\item {\small{}Inference without parametric assumptions (usually) requires
an asymptotic approximation}{\small\par}
\end{itemize}
\end{frame}

\begin{frame}{OLS Asymptotics}

\begin{itemize}
\item {\small{}Let's again express the OLS estimator as}{\small\par}
\end{itemize}
\begin{center}
{\small{}$\hat{\beta}=\beta+\left(\tfrac{1}{N}\sum_{i}X_{i}X_{i}^{\prime}\right)^{-1}\left(\tfrac{1}{N}\sum_{i}X_{i}U_{i}\right)$}{\small\par}
\par\end{center}
\begin{itemize}
\item {\small{}This can be rewritten}{\small\par}
\end{itemize}
\begin{center}
{\small{}$\sqrt{N}\left(\hat{\beta}-\beta\right)=\left(\frac{1}{N}\sum_{i}X_{i}X_{i}^{\prime}\right)^{-1}\left(\tfrac{1}{\sqrt{N}}\sum_{i}X_{i}U_{i}\right)$}{\small\par}
\par\end{center}
\begin{itemize}
\item {\small{}As $N\rightarrow\infty$ this quantity has the same distribution
as:}{\small\par}
\end{itemize}
\begin{center}
{\small{}$E\left[X_{i}X_{i}^{\prime}\right]^{-1}\left(\tfrac{1}{\sqrt{N}}\sum_{i}X_{i}U_{i}\right)$}{\small\par}
\par\end{center}
\begin{itemize}
\item Note: all asymptotic variability through the ``score'' $\sum_{i}X_{i}U_{i}$.
(Why?)
\end{itemize}
\end{frame}

\begin{frame}{OLS Asymptotics}

\begin{itemize}
\item We know $E\left[X_{i}U_{i}\right]=0$, so
\end{itemize}
\begin{center}
$\tfrac{1}{\sqrt{N}}\sum_{i}X_{i}U_{i}\overset{d}{\longrightarrow}N\left(0,E\left[X_{i}X_{i}^{\prime}U_{i}^{2}\right]\right)$
\par\end{center}
\begin{itemize}
\item Therefore
\[
\sqrt{N}\left(\hat{\beta}-\beta\right)\overset{d}{\longrightarrow}N\left(0,E\left[X_{i}X_{i}^{\prime}\right]^{-1}E\left[X_{i}X_{i}^{\prime}U_{i}^{2}\right]E\left[X_{i}X_{i}^{\prime}\right]^{-1}\right)
\]
\item This is the (first-order) asymptotic distribution of the OLS estimator
with \emph{iid} sampling.
\end{itemize}
\end{frame}

\begin{frame}{Heteroskedasticity-robust std errors}

\begin{center}
{\footnotesize{}$\sqrt{N}\left(\hat{\beta}-\beta\right)\overset{d}{\longrightarrow}N\left(0,E\left[X_{i}X_{i}^{\prime}\right]^{-1}E\left[X_{i}X_{i}^{\prime}U_{i}^{2}\right]E\left[X_{i}X_{i}^{\prime}\right]^{-1}\right)$}{\footnotesize\par}
\par\end{center}
\begin{itemize}
\item {\footnotesize{}Under homoskedasticity $E\left[U_{i}^{2}|X_{i}\right]=\sigma^{2}$
and variance collapses to $\sigma^{2}E\left[X_{i}X_{i}^{\prime}\right]^{-1}$}{\small{}\medskip{}
}{\small\par}
\item {\footnotesize{}Modern approach: avoid imposing assumptions on $E\left[U_{i}^{2}|X_{i}\right]$
and base inference on}{\footnotesize\par}
\end{itemize}
\begin{center}
{\footnotesize{}$\hat{V}\left(\hat{\beta}\right)=\frac{1}{N}\left(\frac{1}{N}\sum_{i}X_{i}X_{i}^{\prime}\right)^{-1}\tfrac{1}{N}\sum_{i}X_{i}X_{i}^{\prime}\hat{U}_{i}^{2}\left(\frac{1}{N}\sum_{i}X_{i}X_{i}^{\prime}\right)^{-1}$}{\footnotesize\par}
\par\end{center}
\begin{itemize}
\item {\footnotesize{}This is the White (1980) heteroskedasticity-robust
covariance matrix}{\small{}\medskip{}
}{\small\par}
\item {\footnotesize{}We assumed }\emph{\footnotesize{}iid}{\footnotesize{}
sampling, what if we relax this assumption?}{\footnotesize\par}
\end{itemize}
\end{frame}
\section{Clustering}
\begin{frame}[plain]{Clustering SEs}
\begin{itemize}
\item Traditionally assumed: observations are independent draws from a common distribution.
\medskip
\item Often implausible in serious economic research $\rightarrow$ dependent data (time series?!?).
\medskip
\item We now review how dependent data can affect inference of OLS (easy to generalize to other estimators, see ps)
\end{itemize}
\end{frame}
%%

\begin{frame}{Grouped Data}
\begin{itemize}
\item {\small{}Suppose the data have a group structure, with $J$ groups
and $M_{j}$ individual observations in ``cluster'' $j$, for a
total sample of $N=\sum_{j}M_{j}$. \medskip{}
}{\small\par}
\item {\small{}Regression model is}{\small\par}
\begin{center}
{\small{}$Y_{ij}=X_{ij}^{\prime}\beta+U_{ij}$}{\small\par}
\par\end{center}
\medskip
\item We are worried that $Cov\left(U_{ij},U_{i'j}\right)\neq0$.
\end{itemize}
\end{frame}

\begin{frame}[shrink=10]
\frametitle{Clustering}
\begin{itemize}
\item {\small{}Recall that the OLS variance matrix is given by }{\small\par}
\end{itemize}
\begin{center}
{\small{}$V=\left(X^{\prime}X\right)^{-1}X^{\prime}\Omega X\left(X^{\prime}X\right)^{-1}$}{\small\par}
\par\end{center}
\begin{itemize}
\item {\small{}With correlation within but not between clusters of $U_{ij}$ implies that the error
variance matrix is block diagonal:}{\small\par}
\end{itemize}
\begin{center}
{\small{}$\E[{U}{U'}]=\Omega=\left[\begin{array}{cccc}
\Omega_{1} & 0 & 0 & 0\\
0 & \Omega_{2} & 0 & 0\\
0 & 0 & \ddots & 0\\
0 & 0 & 0 & \Omega_{J}
\end{array}\right]$}{\small\par}
\par\end{center}
\begin{itemize}
\item {\small{}Here $\Omega_{j}$ is the $M_{j}\times M_{j}$ covariance
matrix for cluster $j$\medskip{}
}{\small\par}
\item {\small{}Can't get accurate estimate of each $\Omega_{j}$. But to
estimate the variance of OLS we only need the ``meat'' in the covariance
sandwich:}{\small\par}
\end{itemize}
\begin{center}
{\small{}$X^{\prime}\Omega X=\sum_{j}X_{j}^{\prime}\Omega_{j}X_{j}$}{\small\par}
\par\end{center}
\begin{itemize}
\item Here $X_{j}$ is a $M_{j}\times K$ matrix that stacks the Xs of all the observations belonging to cluster $j$.
\item {\small{}Obvious solution: treat the cluster level data vectors $(X_{j},Y_{j})$
as primitives and apply usual White logic that we have covered before }{\small\par}
\end{itemize}
\end{frame}

\begin{frame}[shrink=10]
\frametitle{Clustering}
\begin{itemize}
\item {\small{}Consistent estimate of the meat in the sandwich can be obtained by invoking a weak LLN for non-identically distributed data (Theorem 3.7 of White, 2001)}{\small\par}
\end{itemize}
\begin{center}
{\small{}$\frac{1}{N}\sum_{j}X_{j}^{\prime}\hat{U}_{j}\hat{U}_{j}^{\prime}X_{j}\overset{p}{\longrightarrow}{\displaystyle \lim_{J\rightarrow\infty}\frac{1}{N}\sum_{j}E\left[X_{j}^{\prime}\Omega_{j}X_{j}\right]}$}{\small\par}
\par\end{center}
\begin{itemize}
\item Recall that in OLS all randomess comes from the score $X'U=\sum_{j=1}^{J}\sum_{i=1}^{M_{j}}X_{ij}U_{ij}$. 
\medskip
\item When normalized by $\frac{1}{N}$ the grand mean  $\frac{1}{N}X'U$ is equivalent to a weighted average of cluster-specific means where the weights crucially depends on $M_{j}$. 
\end{itemize}
\pause
\begin{itemize}
\item {\small{}This generates the ``cluster-robust'' covariance matrix}{\small\par}
\end{itemize}
\begin{center}
{\small{}$\hat{V}\left(\hat{\beta}\right)=\frac{1}{N}\left(\frac{1}{N}\sum_{j}X_{j}^{\prime}X_{j}\right)^{-1}\left(\frac{1}{N}\sum_{j}X_{j}^{\prime}\hat{U}_{j}\hat{U}_{j}^{\prime}X_{j}\right)\left(\frac{1}{N}\sum_{j}X_{j}^{\prime}X_{j}\right)^{-1}$}{\small\par}
\par\end{center}
\begin{itemize}
\item {\small{}This approach relies on asymptotics in the }\emph{\small{}number
of groups; }{\small{}will perform poorly with few clusters\medskip{}
}{\small\par}
\item {\small{}Angrist and Pischke (2009): magic \# of clusters is 42!\medskip{}
}{\small\par}
\end{itemize}
\end{frame}
%%
\begin{frame}[plain]{Comments}

\begin{itemize}
\item As usual, White SEs are biased in finite sample, this is why Stata implements a degrees of freedom correction to regression residuals 
\be
\tilde{{U}}_{j}=\sqrt{\dfrac{J}{J-1}\dfrac{N-1}{N-K}}\hat{{U}}_{j}
\ee
and replaces $\hat{{U}}_{j}$ with $\tilde{{U}}_{j}$ when computing  $\hat{V}\left(\hat{\beta}\right)$.
\pause
\medskip
\item Key question: at what level to cluster the standard error?
\medskip
\begin{itemize}
\item Important: being worried about $Cov\left(U_{ij},U_{i'j}\right)\neq0$ does not \textit{necessarily} means that you have to cluster the standard errors. 
\medskip
\item What one should be worried about is whether $Cov\left(X_{ij}U_{ij},X_{i'j}U_{i'j}\right)\neq0$.

\end{itemize}
\pause
\item Rule of thumb: cluster at the level at which your treatment varies.
\medskip
\end{itemize}
\end{frame}
%%%
\begin{frame}
\frametitle{Example 1: Individual level RCT}
\begin{itemize} 
\item Example: you randomize cash transfers to individuals living in different villages in India.
\medskip
\item In this application: $i$ indexes individuals, $j$ indexes clusters. 
\medskip
\item You want to run the following regression
\medskip
\begin{equation}
    Y_{i} = \alpha + \beta X_{ij} + e_{ij}
\end{equation}
\medskip 
\item Note that here if the treatment variable is randomized at the individual level $ \rightarrow$$\E[X_{ij}X_{i'j}]=0$ (suppose the $X$ variable has been demeaned).
\item This in turn implies that (by LIE) $\E[X_{ij}U_{ij}X_{i'j}U_{i'j}]=\E[X_{ij}X_{i'j}\E[U_{ij}U_{i'j}\vert X_{ij} X_{i'j}]]=0$ $\rightarrow$ no need to cluster!
\end{itemize}
\end{frame}
\begin{frame}
\frametitle{Example 2: Clustered RCT}
\begin{itemize} 
\item Example: you randomize cash transfers to different villages in India.
\medskip
\item You still have data on individuals as well as villages and are interested in running 
\begin{equation}
    Y_{i} = \alpha + \beta X_{ij} + e_{ij}
\end{equation}
\medskip 
\item Now $\E[X_{ij}X_{i'j}]\neq0$ (treatment is perfectly correlated within clusters!). 
\medskip
\item This in turn implies that $\E[X_{ij}U_{ij}X_{i'j}U_{i'j}]=\E[X_{ij}X_{i'j}\E[U_{ij}U_{i'j}\vert X_{ij} X_{i'j}]]\neq0$ $\rightarrow$ now you have to cluster the SEs (at the village level) 
\end{itemize}
\end{frame}
\begin{frame}
\frametitle{Example 3: Real Life}
\begin{itemize} 
\item Most cases treatment of interest might vary within clusters.
\medskip
\item Example: you have data on individuals observed within each state across different years.
\medskip
\item Treatment is a policy implemented by a given state in a particular year (and remains active thereafter).
\medskip
\item Treatment therefore varies at the state-by-year level. 
\medskip
\item However, since treatment is persistent, people often cluster at the level of the State, although such SEs might be too conservative (Bertrand, Duflo and Mullainathan, 2004; Abadie et al., 2023).
\end{itemize}
\end{frame}

%%
\begin{frame}[plain]{Key Reference}

\begin{center}
\includegraphics[scale=0.2]{aux/abadie_QJE}
\par\end{center}

\end{frame}

\begin{frame}[plain]{Cluster SEs in Action}

\begin{center}
\includegraphics[scale=0.2]{aux/abadie22.jpg}
\par\end{center}
\end{frame}%%

\begin{frame}{Takeaways from AAIW (2023)}
\begin{itemize}
\item Figure it out at what level to cluster requires figuring out (i) the sampling process and (ii) how is treatment assigned. 
\medskip
\item Thinking of a framework where $J\rightarrow\infty$ is often not useful $\rightarrow$ AAIW shows that is possible to derive standard errors smaller than those based on asymptotic approximation that refers to infinite super-population.
\medskip
\item Correlation in the regression error term within clusters does not automatically mean that one has to cluster (see example \#1).
\medskip
\item The idea of  "keep clustering at higher levels if SEs continue to become smaller" is ill-advised. 
\medskip
\item Derive two new cluster-robust estimator of SEs that are considerably smaller than White-based cluster-robust SEs.
\begin{itemize}
    \item Especially valuable in contexts like Example 3 (``partially" clustered treatment assignment). 
\end{itemize}
\end{itemize}
    
\end{frame}

\begin{frame}[plain]{The need for clustering...}

\begin{center}
\includegraphics[scale=0.1]{aux/clustering}
\par\end{center}
\end{frame}%%

\section{Bootstrap}
\begin{frame}[plain]{The Bootstrap}

\begin{center}
\includegraphics[scale=0.3]{aux/space_jam}
\par\end{center}

\end{frame}

\begin{frame}[plain]{Preliminaries}

\begin{itemize}
\item A test statistic can be written: 
\[
T_{n}\equiv T_{n}\left(\boldsymbol{z}_{1},\boldsymbol{z}_{2},...,\boldsymbol{z}_{N},F\right)
\]
where $F$ is population distribution of data
\item Example: ``t-stat'' is $T_{n}=\frac{\hat{\theta}-\theta_{0}}{\hat{\sigma}}$
\end{itemize}
\medskip{}

\begin{itemize}
\item Finite sample dist of $T_{n}$ is 
\[
G_{n}\left(t,F\right)\equiv P\left(T_{n}\leq t|F\right)
\]
\item Problem: don't know $F$! (except under very strong null hypothesis)
\end{itemize}
\end{frame}

\begin{frame}[plain]{Two approaches}

Conventional asymptotics:
\[
G_{n}\left(t,F\right)\approx\underset{n\rightarrow\infty}{\lim}G_{n}\left(t,F\right)=G\left(t\right)\mbox{ ("asymptotic pivot")}
\]

\begin{itemize}
\item Example: ``t-stat'' goes to normal $G\left(t\right)=\Phi\left(t\right)$
\end{itemize}

\pause{\medskip{}
}

Bootstrap approximation:
\[
G_{n}\left(t,F\right)\approx G_{n}\left(t,F_{n}\right)
\]

\begin{itemize}
\item Traditional choice of $F_{n}$ is EDF
\item Glivenko-Cantelli: EDF converges uniformly to $F$
\end{itemize}
\end{frame}

\begin{frame}[plain]{Convergence in action}

\begin{center}
\includegraphics[scale=0.3]{aux/glivenko}
\par\end{center}

Source: Hansen online textbook 
\end{frame}

\begin{frame}[plain]{Computation (Pairs Bootrstrap)}

For small $N$, possible to enumerate $G_{n}\left(t,F_{n}\right)$
exactly ($2^{N}$ atoms)

Otherwise use sampling w/ replacement to approximate $G_{n}\left(t,F_{n}\right)$

\pause{\medskip{}
}

Algorithm:
\begin{enumerate}
\item Draw a random sample of size $N$ \textit{with replacement} from the
data. Call this the $b$'th ``bootstrap\textquotedblright{} sample.
\item Compute the test statistic $T_{n}^{\ast(b)}$ in the bootstrap sample
(see below for details on this).
\item Repeat \#1 and \#2 until $b=B$.
\end{enumerate}

\pause{\medskip{}
}

Some rules of thumb for $B$ (from my personal experience):
\begin{itemize}
\item $B\geq400$ for std error estimation
\item $B\geq999$ for quantiles
\end{itemize}
\medskip
What if errors are serially correlated within a cluster? \pause Block Bootstrap.
What if I can't compute my model in a given bootrstrap replica? \pause Residual/Wild Bootstrap.
\end{frame}

\begin{frame}[plain]{When does the bootstrap work?}

\begin{itemize}
\item Short answer: when things are ``well behaved''
\item Problems emerge when limiting distribution is not normal
\item Example (Andrews, 2000):
\[
\theta_{0}=\max\left\{ 0,E\left[X_{i}\right]\right\} 
\]
\[
T_{n}=\sqrt{N}\left(\max\left\{ 0,\overline{X}\right\} -\theta_{0}\right)
\]
\item $T_{n}$ is not ``regular'' -- limiting distribution is normal
when $\theta_{0}>0$ but ``half-normal'' when $\theta_{0}=0$. 
\item Discontinuity at 0 creates problems for bootstrap
\item Practical areas where this problem arises: 

\begin{itemize}
\item Partially identified models (Canay and Shaikh, 2016)
\item Nearest neighbor matching (Abadie and Imbens, 2008)
\end{itemize}
\end{itemize}
\end{frame}

\begin{frame}[plain]{A useful guide..}

\begin{theorem}[Fang and Santos, 2015]
 If $\sqrt{N}\left(\hat{\theta}-\theta_{0}\right)\overset{d}{\rightarrow}N\left(0,\sigma\right)$
and the bootstrap estimate $\theta^{*}$ obeys $\sqrt{N}\left(\theta^{*}-\hat{\theta}\right)\overset{d}{\rightarrow}N\left(0,\sigma\right)$,
then the bootstrap is consistent for the limiting distribution of
$\sqrt{N}\left(f\left(\hat{\theta}\right)-f\left(\theta_{0}\right)\right)$
\emph{iff} $f\left(.\right)$ is differentiable.
\end{theorem}
\medskip{}

\begin{itemize}
\item ``Delta method'' for the bootstrap
\item Rules out $\max\left\{ .\right\} /\min\left\{ .\right\} $ transforms
and ``kinks''
\item But allows smooth functions of sample moments.
\end{itemize}
\end{frame}

\begin{frame}[plain]{When to use the bootstrap}

\begin{itemize}
\item When computing an analytical standard error is tedious:

\begin{itemize}
\item Quantile regression (avoid estimation of nuisance density)
\item Multi-step estimation (just bootstrap the entire procedure)
\end{itemize}
\end{itemize}
\bigskip{}

\begin{itemize}
\item When tests based on analytical standard error are poorly behaved (though
other options are typically available):

\begin{itemize}
\item Few clusters
\item Delta method for highly nonlinear hypothesis
\item Weak IV
\end{itemize}
\end{itemize}
\end{frame}

\begin{frame}[plain]{Bootstrap refinements}

\begin{itemize}
\item Bootstrapping t-stat often better than comparing t-stat to normal
crit value (i.e. 1.96)
\item Provides higher order approximation to $G_{n}\left(t,F\right)$ when
$T_{n}\left(\theta\right)\equiv\frac{\widehat{\theta}-\theta}{\widehat{\sigma}}$
is an asymptotic pivot (Hall, 1992)

\begin{itemize}
\item note: requires consistent std err estimator $\hat{\sigma}$!
\end{itemize}
\medskip{}

\item Algorithm:
\end{itemize}
\begin{enumerate}
\item Draw $B$ bootstrap samples
\item Form BS test statistics $T_{n}^{\ast\left(b\right)}\left(\theta\right)\equiv\frac{\widehat{\theta}^{\ast\left(b\right)}-\theta}{\widehat{\sigma}^{\ast\left(b\right)}}$
\item Find $1-\alpha$'th quantile of $\left\vert T_{n}^{\ast\left(b\right)}\left(\hat{\theta}\right)\right\vert $
, call it $c^{*}$
\item Reject if $\left\vert T_{n}\left(\theta_{0}\right)\right\vert >c^{*}$
\end{enumerate}
\end{frame}

\begin{frame}[plain]{Kline and Walters (2016, QJE)}

\includegraphics[scale=0.5]{aux/kline_walters_QJE}
\begin{itemize}
\item t-stat for $MVPF=1$ in ``pessimistic'' scenario is $\frac{.5}{.34}=1.47$
\item Normal p-val for 1-sided test would be $1-\Phi\left(1.47\right)=.07$.
Refined p-val$\approx0$!
\end{itemize}
\end{frame}

\begin{frame}[plain]{A monte carlo study}

\begin{center}
\includegraphics[scale=0.2]{aux/kline_walters_MC}
\par\end{center}
\begin{itemize}
\item Delta method normal approximation is poor here

\begin{itemize}
\item Overcoverage: delta method rejected in 0/500 simulations
\end{itemize}
\item Bootstrap picks up asymmetry in finite sample dist

\begin{itemize}
\item Good coverage: 7\% rejections on nominal 5\% test
\end{itemize}
\end{itemize}
\end{frame}

\begin{frame}[plain]{Residual Bootstraps}

Residual bootstrap for linear model (Freedman, 1981)

Algorithm:
\begin{enumerate}
\item Fit linear model $Y_{i}=X_{i}\beta+\varepsilon_{i}$
\item Resample from residuals $\hat{\varepsilon}_{i}$ to make new dependent
variable $Y_{i}^{*}=X_{i}\hat{\beta}+\varepsilon_{i}^{*}$
\item Fit linear model to $Y_{i}^{*}$, get $\beta^{*}$
\end{enumerate}

\pause{\medskip{}
}
\begin{itemize}
\item Very good asymptotic properties

\begin{itemize}
\item Can even allow $dim\left(\beta\right)$ to grow w/ $N$ (Bickel and
Freedman, 1983)
\item Intuition: resampling from a scalar instead of from $\left(Y_{i},X_{i}\right)$
\end{itemize}
\item Do you see a problem with this?
\end{itemize}
\end{frame}

\begin{frame}[plain]{Wild Bootstrap (Wu, 1986; Liu, 1988):}

Allow heteroscedasticity while retaining advantage of residual BS

Algorithm:
\begin{enumerate}
\item Fit linear model $Y_{i}=X_{i}\beta+\varepsilon_{i}$
\item Draw iid weights $\left\{ W_{i}^{*}\right\} _{i=1}^{N}$ obeying $E\left[W_{i}^{*}\right]=0$,
$E\left[\left(W_{i}^{*}\right)^{2}\right]=1$ 
\item Make new dependent variable $Y_{i}^{*}=X_{i}\hat{\beta}+\hat{\varepsilon}_{i}W_{i}^{*}$
\item Fit linear model to $Y_{i}^{*}$, get $\left(\beta^{*},\sigma^{*}\right)$
\end{enumerate}

\pause{}

Intuition:
\[
E\left[\hat{\varepsilon}_{i}W_{i}^{*}|\left\{ Y_{i},X_{i}\right\} _{i=1}^{N}\right]=0
\]
\[
E\left[\left(\hat{\varepsilon}_{i}W_{i}^{*}\right)^{2}|\left\{ Y_{i},X_{i}\right\} _{i=1}^{N}\right]=\hat{\varepsilon}_{i}^{2}
\]

So $\varepsilon_{i}^{*}\equiv\hat{\varepsilon}_{i}W_{i}^{*}$ is mean
independent of $X_{i}$ and replicates hetero found in data

But do we want to impose mean independence?
\end{frame}

\begin{frame}[plain]{Theoretical results}

\begin{itemize}
\item Mammen (1993): 

\begin{itemize}
\item Wild BS yields refinement even under many regressors asymptotics $\left(p\rightarrow\infty,\mbox{ }p/n\rightarrow0\right)$
provided linear model is correct
\item Wild (\& ``pairs'') BS still consistent under misspecification
\end{itemize}
\medskip{}

\item Kline and Santos (2012, JoE): Wild BS retains refinement under fixed
$p$ asymptotics when ``locally'' misspecified
\end{itemize}
\medskip{}

\begin{itemize}
\item Kline and Santos (2012, JEMS): ``Score BS''

\begin{itemize}
\item extension of Wild BS to nonlinear models (e.g. Probit)
\item avoid recomputation of estimator (major speedup)
\item connection to Bayesian/Weighted BS
\item implemented in ``\href{https://ideas.repec.org/c/boc/bocode/s458121.html}{BOOTTEST}''
Stata package
\end{itemize}
\end{itemize}
\end{frame}

\begin{frame}[plain]{Cameron, Gelbach, and Miller (2008, RESTAT)}

\begin{itemize}
\item Monte Carlo bootstrap tests in clustered data setting
\item Variant of Wild BS that ``imposes the null'' performs best
\end{itemize}
Algorithm:
\begin{enumerate}
\item Fit \emph{restricted} linear model $Y_{i}=X_{i}\tilde{\beta}+\varepsilon_{i}$
where $c\left(\tilde{\beta}\right)=0$
\item Draw iid weights $\left\{ W_{i}^{*}\right\} _{i=1}^{N}$ obeying $E\left[W_{i}^{*}\right]=0$,
$E\left[\left(W_{i}^{*}\right)^{2}\right]=1$ 
\item Make new dependent variable $Y_{i}^{*}=X_{i}\hat{\tilde{\beta}}+\hat{\tilde{\varepsilon}}_{i}W_{i}^{*}$
\item Fit \emph{unrestricted }linear model to $Y_{i}^{*}$, get $\left(\beta^{*},\sigma^{*}\right)$
\item Compare clustered Wald stat $\hat{T}_{n}$ to appropriate quantile
of bootstrap analogue $T_{n}^{*}$
\end{enumerate}

\pause{\smallskip{}
}
\begin{itemize}
\item Canay, Shaikh, Santos (2016): CGM BS consistent for fixed \# of clusters
as cluster size grows, given balanced design ($E\left[X_{i}X_{i}'|C_{i}\right]=E\left[X_{i}X_{i}'\right]$)
and effect homogeneity $\beta_{c}=\beta$.
\end{itemize}
\end{frame}
%%
\begin{frame}
\frametitle{Beware!}
\includegraphics[scale=0.4]{aux/bootstrap_high_dimensions}
\end{frame}
%%
\begin{frame}
\frametitle{(The answer is no)}
\includegraphics[scale=0.4]{aux/bootstrap_high_dimensions}
\end{frame}

\section{Two-Step Estimators}

\begin{frame}[plain]{2-step estimators}

Lots of settings where it is easier to estimate model in ``steps''
\begin{itemize}
\item 1st step parameter estimate $\hat{\gamma}$
\item Second step m-estimator:
\[
\widehat{\theta}\equiv\underset{\theta\in\Theta}{\arg\min}\sum\limits _{i}q\left(\boldsymbol{Z}_{i},\theta;\widehat{\gamma}\right)
\]
\end{itemize}
Example (WNLLS):
\[
\widehat{\theta}=\underset{\theta\in\Theta}{\arg\min}\sum_{i}\left[Y_{i}-m\left(X_{i},\theta\right)\right]^{2}w\left(X_{i};\widehat{\gamma}\right)
\]

\begin{itemize}
\item $\hat{\gamma}\pto\gamma^{*}$
\item How/when to adjust standard errors for use of $\hat{\gamma}$ instead
of $\gamma^{*}$?
\end{itemize}
\end{frame}

\begin{frame}[plain]{Influence fn representation}

Asymptotic variability derives from two sources:
\begin{eqnarray*}
\sqrt{N}\left(\widehat{\theta}-\theta_{0}\right) & = & -E\left[\boldsymbol{H}_{i}\left(\theta_{0},\gamma^{\ast}\right)\right]^{-1}N^{-1/2}\sum_{i}\boldsymbol{s}_{i}\left(\theta_{0},\gamma^{\ast}\right)\\
 &  & -E\left[\boldsymbol{H}_{i}\left(\theta_{0},\gamma^{\ast}\right)\right]^{-1}E\left[\bigtriangledown_{\gamma}\boldsymbol{s}_{i}\left(\theta_{0},\gamma^{\ast}\right)\right]\sqrt{N}\left(\widehat{\gamma}-\gamma^{\ast}\right)\\
 &  & +o_{p}\left(1\right)
\end{eqnarray*}
Second term is zero when $E\left[\bigtriangledown_{\gamma}\boldsymbol{s}_{i}\left(\theta_{0},\gamma^{\ast}\right)\right]=0$!

\medskip{}

Example (WNLLS):
\[
\bigtriangledown_{\gamma}\boldsymbol{s}_{i}\left(\theta_{0},\gamma^{\ast}\right)=\left[Y_{i}-m\left(X_{i},\theta_{0}\right)\right]\frac{\partial}{\partial\theta}m\left(X_{i},\theta_{0}\right)\bigtriangledown_{\gamma}w\left(X_{i};\gamma^{\ast}\right)
\]

\begin{itemize}
\item If $E\left[Y_{i}|X_{i}\right]=m\left(x_{i},\theta_{0}\right)$ then
$E\left[\bigtriangledown_{\gamma}\boldsymbol{s}_{i}\left(\theta_{0},\gamma^{\ast}\right)\right]=0$
(by LIE)
\item Otherwise, all bets are off
\end{itemize}
\end{frame}
%%
\begin{frame}[plain]{Two-Steps estimators are everywhere}


\begin{center}
\includegraphics[scale=0.3]{aux/chetty3}
\par\end{center}


\end{frame}



\begin{frame}[plain]{Place ``effects" }


\begin{center}
\includegraphics[scale=0.45]{aux/chetty1}
\par\end{center}


\end{frame}

\begin{frame}[plain]{Two-step in action }


\begin{center}
\includegraphics[scale=0.45]{aux/chetty2}
\par\end{center}


\end{frame}
\end{document}
